\section{DeepSeek MoE case study：MoE、MLA、MTP 的组合}

本节把 DeepSeek v1/v2/v3 作为案例，观察 shared experts、fine-grained experts、MLA 和 MTP 如何组合。


\begin{knowledgebox}{first-use glossary：MLA 与 MTP}
MLA 是 Multi-head Latent Attention，把 K/V cache 压缩为低维 latent 后再投影恢复，降低长上下文推理内存。MTP 是 Multi-Token Prediction，让轻量模块预测多个未来 token，可用于改善训练信号或推理草稿。DeepSeek v3 把 MoE、MLA、MTP 等选择组合成一套系统 recipe。
\end{knowledgebox}


\begin{figure}[H]
\centering
\includegraphics[width=0.88\textwidth]{slide-053.jpg}
\caption{Slide 54：DeepSeek MoE case study 的路线图。}
\figsource{CS336 Spring 2026 Lecture 4 官方幻灯片。}
\end{figure}

\noindent\textbf{展开说明：}To wrap up, we’ll walk through the DeepSeek MoE architecture. 这组页面不是模型宣传，而是检查前文机制如何组合：fine-grained/shared experts 改变容量，auxiliary-loss-free balancing 改变路由稳定性，MLA 压缩 inference state，MTP 增加训练信号并可能支持 speculative decoding。

\begin{knowledgebox}{讲义补充：源 Slide 54 应该怎么看}
这页是机制或证据页。先看它比较的是 compute、参数量、routing、负载均衡还是通信；再看结论支持的是质量提升、训练速度、推理成本还是系统复杂度。MoE 和 attention alternatives 的图常把模型结构和系统代价混在一起，读图时要分清“算法机制”和“硬件执行成本”。
\end{knowledgebox}



\begin{figure}[H]
\centering
\includegraphics[width=0.88\textwidth]{slide-054.jpg}
\caption{Slide 55：DeepSeek MoE v2 的总参数、active parameters 与架构组件。}
\figsource{CS336 Spring 2026 Lecture 4 官方幻灯片。}
\end{figure}

\noindent\textbf{展开说明：}V2 (236B – 21 active): 总参数和 active parameters 必须分开读。前者说明模型能容纳的专家容量，后者近似每个 token 的计算；但服务仍需存放或跨节点访问完整专家集合，因此不能用 21B active 把它当成普通 21B dense model。

\begin{knowledgebox}{讲义补充：源 Slide 55 应该怎么看}
这页是机制或证据页。先看它比较的是 compute、参数量、routing、负载均衡还是通信；再看结论支持的是质量提升、训练速度、推理成本还是系统复杂度。MoE 和 attention alternatives 的图常把模型结构和系统代价混在一起，读图时要分清“算法机制”和“硬件执行成本”。
\end{knowledgebox}


\begin{figure}[H]
\centering
\includegraphics[width=0.88\textwidth]{slide-055.jpg}
\caption{Slide 56: DeepSeek MoE v3}
\figsource{CS336 Spring 2026 Lecture 4 官方幻灯片。}
\end{figure}

\noindent\textbf{展开说明：}V2 (671B – 37 active):

\begin{knowledgebox}{读图：Slide 56 应该怎么看}
这页是机制或证据页。先看它比较的是 compute、参数量、routing、负载均衡还是通信；再看结论支持的是质量提升、训练速度、推理成本还是系统复杂度。MoE 和 attention alternatives 的图常把模型结构和系统代价混在一起，读图时要分清“算法机制”和“硬件执行成本”。
\end{knowledgebox}



\begin{figure}[H]
\centering
\includegraphics[width=0.88\textwidth]{slide-056.jpg}
\caption{Slide 57：DeepSeek v3 attention 中 MLA 与 latent state 的位置。}
\figsource{CS336 Spring 2026 Lecture 4 官方幻灯片。}
\end{figure}

\noindent\textbf{展开说明：}MLA : Multihead latent attention。它把可缓存的 K/V 信息先压到低维 latent，再通过投影参与 attention，从而减少长上下文 decode 的 cache bytes；收益来自 memory traffic，而不是让 attention 数学本身消失。

\begin{knowledgebox}{讲义补充：源 Slide 57 应该怎么看}
这页是机制或证据页。先看它比较的是 compute、参数量、routing、负载均衡还是通信；再看结论支持的是质量提升、训练速度、推理成本还是系统复杂度。MoE 和 attention alternatives 的图常把模型结构和系统代价混在一起，读图时要分清“算法机制”和“硬件执行成本”。
\end{knowledgebox}



\begin{figure}[H]
\centering
\includegraphics[width=0.88\textwidth]{slide-057.jpg}
\caption{Slide 58：MLA 的低维 latent 分解与 Q/K/V 重构。}
\figsource{CS336 Spring 2026 Lecture 4 官方幻灯片。}
\end{figure}

\noindent\textbf{展开说明：}Basic idea: express the Q, K, V as functions of a lower-dim latent activation。读公式时先区分训练投影 FLOPs 与推理缓存大小：低秩表示可能增加投影步骤，却能显著减少每层、每 token 需要长期保存和读取的状态。

\begin{importantbox}{DeepSeek case study 的组合逻辑}
MoE、MLA 与 MTP 解决的资源不同：MoE 控制每 token 激活参数，MLA 控制 decode KV/state memory，MTP 增加多步预测监督并为草稿生成提供接口。把它们放在一起时必须同时检查 router balance、expert communication、latent projection cost、cache bytes 和多 token 目标权重；任何单项收益都可能被另一项的系统开销抵消。
\end{importantbox}

\begin{knowledgebox}{讲义补充：源 Slide 58 应该怎么看}
这页是机制或证据页。先看它比较的是 compute、参数量、routing、负载均衡还是通信；再看结论支持的是质量提升、训练速度、推理成本还是系统复杂度。MoE 和 attention alternatives 的图常把模型结构和系统代价混在一起，读图时要分清“算法机制”和“硬件执行成本”。
\end{knowledgebox}


\begin{figure}[H]
\centering
\includegraphics[width=0.88\textwidth]{slide-058.jpg}
\caption{Slide 59: What else do you need to make DeepSeek MoE v3?}
\figsource{CS336 Spring 2026 Lecture 4 官方幻灯片。}
\end{figure}

\noindent\textbf{展开说明：}MTP: Have small, lightweight models that predict multiple steps ahead

\begin{knowledgebox}{读图：Slide 59 应该怎么看}
这页是机制或证据页。先看它比较的是 compute、参数量、routing、负载均衡还是通信；再看结论支持的是质量提升、训练速度、推理成本还是系统复杂度。MoE 和 attention alternatives 的图常把模型结构和系统代价混在一起，读图时要分清“算法机制”和“硬件执行成本”。
\end{knowledgebox}

\subsection{本章小结}
本节的共同线索是 selective computation：通过结构选择让 token 不必总是访问所有历史或所有参数。但选择性越强，越需要额外机制保证信息流、负载均衡和训练稳定。

\subsection{Selective computation 的资源验收表}

现在所有教学 slide 都已按源顺序纳入，最后需要把图中的机制转成可执行验收。Attention alternatives 与 MoE 都会把一类 dense cost 变成新的稀疏状态和调度成本；如果只测理论复杂度，很容易得到方向正确但系统错误的结论。下面的账本用于小规模原型、scaling experiment 和服务 benchmark：每一行都应有实际 profiler 或日志，而不是只写架构名称。

\begin{longtable}{p{0.20\textwidth}p{0.32\textwidth}p{0.38\textwidth}}
\toprule
机制 & 节省什么 & 必须新增测量什么 \\
\midrule
Linear/recurrent attention & 避免完整 $n\times n$ attention matrix，压缩长序列计算或状态。 & state size、update kernel、长距离任务质量、prefill/decode latency。 \\
Sparse/local attention & 只访问少量历史位置或局部窗口。 & 选中位置数、索引/稀疏 kernel 开销、全局信息恢复周期、长程 retrieval。 \\
Top-$k$ MoE & 每个 token 只激活少数 experts，以较小 active FLOPs 使用更大总容量。 & total/active params、router entropy、expert utilization、capacity overflow。 \\
Expert parallelism & 把完整 expert weights 分散到多设备，扩展总参数。 & dispatch/combine bytes、all-to-all 时间、尾部设备负载、通信计算重叠率。 \\
Upcycling & 复用 dense checkpoint，减少 MoE 从头预训练成本。 & 初始化函数等价性、expert 分化速度、router collapse、继续训练 token 数。 \\
MLA & 压缩长期缓存的 K/V 信息，降低 decode memory traffic。 & latent/cache bytes、重构投影 FLOPs、quality delta、batch/sequence 长度下的吞吐。 \\
MTP & 一次训练多个未来 token，并可提供 speculative draft。 & 辅助 loss 权重、acceptance rate、额外训练 FLOPs、端到端 decode speedup。 \\
\bottomrule
\end{longtable}

\begin{importantbox}{最小可复现实验：不要一上来训练超大 MoE}
先用相同 tokenizer、数据顺序、hidden size 与 token budget 构造三个小模型：一个 dense FFN baseline，一个总参数更大但 active FLOPs 对齐的 top-2 MoE，一个加入 expert parallelism 的同配置 MoE。训练时同时记录 validation loss、router entropy、每个 expert 的 token histogram、capacity overflow、dispatch/combine bytes、expert GEMM 时间和完整 step time。随后把 batch size 或设备数改到未参与调参的 hold-out setting；若 MoE 只在原始设置中更快，说明收益来自偶然的 kernel shape 或通信拓扑，而不是可迁移的 conditional-compute recipe。对 attention alternatives 也采用相同原则：固定模型和数据，只改变 state/window 机制，并在短序列、目标长序列与未见过的更长序列上同时比较质量和 latency。
\end{importantbox}

\begin{warningbox}{稀疏性会把平均成本问题变成尾延迟问题}
Dense layer 的每个 token 工作量接近固定，而 routing 后不同 experts 的 token 数会波动。平均 FLOPs 很低时，最忙 expert、最慢通信 rank 或容量溢出可能决定整个 step；因此必须看分位数、最大负载和跨 rank variance，而不能只报告平均 expert utilization。线上推理还要分别测单请求 latency 与高并发 throughput：较大的 batch 可能改善 expert GEMM 利用率，却也会增加排队、路由聚合和尾部等待。只有在目标并发、上下文长度和硬件拓扑下都复验，才能把实验室稀疏性称为部署收益。训练侧也应保留逐步 expert-load 日志，因为最终平均均衡可能掩盖早期长时间 collapse；这种早期失衡会改变各 expert 接触的数据分布，即使后期 loss 恢复，也不代表训练轨迹等价。
还要结合多个 checkpoint 的路由分布复核。
\end{warningbox}

